% ----------------------------------------------------------------------
%  3.8. Упругие свойства: 3.8.1–3.8.3 (характеристики упругости,
%  упругость металлов и сплавов, ферромагнитная аномалия)
% ----------------------------------------------------------------------
\begin{orientir}
Упругие константы нужны конструктору для расчётов, а физику и
материаловеду~— как чуткий индикатор межатомной связи и фазовых
превращений. Мы введём основные характеристики упругости и способы их
измерения (3.8.1), посмотрим, как модуль зависит от положения элемента в
таблице Менделеева и от состава сплава (3.8.2), разберём аномалию
упругости ферромагнетиков (3.8.3). Затем выйдем за рамки закона Гука:
деформация отстаёт от напряжения~— это релаксация (3.8.4) и внутреннее
трение (3.8.5).
\end{orientir}

\subsection{Основные характеристики упругости}\label{ssec:elastic-const}

Упругость изотропного тела описывают четыре величины:
\begin{itemize}
\item $E$~— \term{модуль нормальной упругости} (модуль Юнга);
\item $G$~— \term{модуль сдвига};
\item $D$~— \term{модуль всестороннего сжатия} (объёмной упругости);
\item $\nu$~— \term{коэффициент Пуассона}.
\end{itemize}
Модули имеют размерность напряжения (обычно ГПа), $\nu$ безразмерен.
Независимы только две из четырёх величин:
\begin{keyeq}\label{eq:G-D-E}
G=\frac{E}{2(1+\nu)},\qquad D=\frac{E}{3(1-2\nu)} .
\end{keyeq}

\paragraph{Закон Гука и коэффициент Пуассона}
При одноосном растяжении тонкого стержня напряжением $\sigma_x$ в упругой
области
\begin{equation}\label{eq:hooke-rod}
\sigma_x=E\varepsilon_x,\qquad \varepsilon_y=\varepsilon_z=-\nu\varepsilon_x ,
\end{equation}
т.\,е. удлинение сопровождается поперечным сужением. Относительное
изменение объёма
$\Delta V/V_0=\varepsilon_x+\varepsilon_y+\varepsilon_z=(1-2\nu)\varepsilon_x$.
Если бы объём при растяжении не менялся, было бы $\nu=0{,}5$; у реальных
твёрдых тел всегда $\nu<0{,}5$. Для большинства металлов и сплавов
$\nu=0{,}25$--$0{,}35$; наименьший коэффициент~— у бериллия
(${\approx}0{,}03$)\ispr{В оригинале для бериллия указано $\nu=0{,}0039$;
справочные значения для поликристаллического бериллия~—
0,02--0,07.}, наибольшие~— у свинца (0,44) и индия (0,46). Поскольку $\nu$
у металлов близки, близки и отношения модулей: при $\nu\approx1/3$ из
\eqref{eq:G-D-E} $G/E\approx3/8$, что и наблюдается.

\paragraph{Анизотропия}
Упругость определяется силами связи между соседними атомами, поэтому
модули монокристалла~— анизотропные величины, зависящие от
направления (тензорное описание~— в п.~1.5.4). В поликристалле из
множества беспорядочно ориентированных зёрен модули от направления не
зависят; их можно рассчитать усреднением констант монокристалла.

\paragraph{Модуль и прочность связи}
Модуль упругости и температура плавления характеризуют одно и то же~—
прочность межатомной связи. Эмпирическое уравнение Портевена:
\begin{equation}\label{eq:portevin}
E=c\,\frac{\Tpl^{\,a}}{V^{\,b}},\qquad a\approx b\approx1 ,
\end{equation}
($V$~— удельный, точнее атомный, объём; $c$~— постоянная); оно хорошо
выполняется в ряду металлов с одинаковой решёткой.

\paragraph{Температурная зависимость модулей}
С нагревом средняя энергия связи уменьшается (решётка расширяется), и
модули падают. Вводят \term{температурный коэффициент модуля}
\begin{keyeq}\label{eq:E-T}
e=-\frac1E\frac{\d E}{\d T},\qquad E\approx E_0(1-eT),\qquad \frac{\alpha}{e}\approx0{,}04 ,
\end{keyeq}
который пропорционален коэффициенту объёмного расширения $\beta=3\alpha$
(это следует из дифференцирования~\eqref{eq:portevin})\ispr{В оригинале
коэффициент записан как $e=(\d E/\d T)\cdot1/T$; правильно
$e=-(\d E/\d T)/E$~— относительное изменение модуля на градус.}. Значит, $e$
ведёт себя как коэффициент расширения, а тот~— как теплоёмкость
(п.~\ref{ssec:expansion}): $e\propto(T/\theta_D)^3$ при $T\ll\theta_D$ и
$e\approx\text{const}$ при $T\gtrsim\theta_D$, т.\,е. выше $\theta_D$ модуль
падает линейно. Вблизи $\Tpl$ спад ускоряется из-за равновесных вакансий.
От низких температур до плавления модули металлов уменьшаются на
несколько десятков процентов. У монокристаллов температурный коэффициент
максимален в направлениях, где модуль минимален, и наоборот.

\paragraph{Модуль, скорость звука и температура Дебая}
Модули определяют скорости упругих волн. В тонком стержне продольная
волна бежит со скоростью $\upsilon_\text{ст}=\sqrt{E/\rho}$, крутильная
(поперечная)~— со скоростью $\upsilon_t=\sqrt{G/\rho}$; в неограниченной
среде скорость продольной волны
$\upsilon_l=\sqrt{(D+\frac43G)/\rho}$. Через среднюю скорость звука
$\upsilon_m$ модули определяют и температуру Дебая (п.~\ref{ssec:heat-capacity}):
\begin{keyeq}\label{eq:thetaD-elastic}
\theta_D=\frac hk\left(\frac{3N_A\rho}{4\pi A}\right)^{1/3}\upsilon_m,\qquad
\frac{3}{\upsilon_m^3}=\frac{1}{\upsilon_l^3}+\frac{2}{\upsilon_t^3},
\end{keyeq}
где $A$~— атомная масса (кг/моль), $\rho$~— плотность\ispr{В оригинале
показатель степени в формуле для $\theta_D$ равен $1/2$; правильно~— $1/3$
(дебаевская частота $\nu_D=\upsilon_m(3n/4\pi)^{1/3}$, $n$~— число атомов в
единице объёма).}. Формула~— прямая связь между «механикой» (модули) и
«термодинамикой» (теплоёмкость).

\begin{primer}[Температура Дебая меди по скорости звука]
Для меди $\rho=8960$~кг/м$^3$, $A=0{,}0635$~кг/моль, $\upsilon_l=4760$~м/с,
$\upsilon_t=2325$~м/с. Средняя скорость:
$3/\upsilon_m^3=1/4760^3+2/2325^3=1{,}68\cdot10^{-10}$~с$^3$/м$^3$,
$\upsilon_m\approx2610$~м/с. Концентрация атомов
$n=N_A\rho/A=8{,}5\cdot10^{28}$~м$^{-3}$, $(3n/4\pi)^{1/3}=2{,}73\cdot10^{9}$~м$^{-1}$.
Тогда $\theta_D=\frac{6{,}63\cdot10^{-34}}{1{,}38\cdot10^{-23}}\cdot2{,}73\cdot10^9\cdot2610\approx342$~К~—
практически совпадает с калориметрическим значением 343~К.
\end{primer}

\paragraph{Динамические методы измерения}
На связи модулей со скоростью звука основаны \emph{динамические} методы.
Они заметно точнее статических (по измеренным напряжению и деформации,
$E=\sigma/\varepsilon$) и позволяют работать при очень малых деформациях,
что важно для мягких металлов, у которых пластическая деформация
начинается уже при малых напряжениях. Методы бывают \emph{импульсные}
(измеряют время пробега импульса деформации через образец) и
\emph{резонансные} (измеряют собственные частоты продольных, изгибных или
крутильных колебаний). Колебания возбуждают и регистрируют
электромагнитными, электростатическими, пьезоэлектрическими или
магнитострикционными преобразователями. Скорость звука в металлах
велика (в свинце около 1300, в меди около 3700, в железе около 5000~м/с),
поэтому работают на звуковых и ультразвуковых частотах.

\paragraph{Резонансный метод}
Стержень длиной $L$ закреплён посередине; на одном конце преобразователь
возбуждает продольные колебания, на другом датчик измеряет их амплитуду.
Когда частота генератора совпадает с собственной частотой стержня,
амплитуда резко возрастает~— резонанс. В точке закрепления~— узел, на
свободных концах~— пучности, поэтому на длине стержня укладывается
нечётное число полуволн: $L=(2n+1)\lambda/2$ (рис.~\ref{fig:rod-res}).
Самый сильный резонанс~— основной, $n=0$, $L=\lambda/2$. Так как
$\lambda=\upsilon_\text{ст}/\nu_\text{рез}$ и $\upsilon_\text{ст}=\sqrt{E/\rho}$,
\begin{keyeq}\label{eq:E-res}
E=4\rho\,\nu_\text{рез}^2L^2 .
\end{keyeq}
Резонансную частоту можно измерить очень точно, отсюда и высокая
точность модуля.

\begin{figure}[ht]
\centering
\begin{tikzpicture}[x=1cm, y=1cm]
  \foreach \yy/\m/\lab in {0/1/{основной резонанс, $L=\lambda/2$}, -2.1/3/{второй резонанс, $L=3\lambda/2$}}{
    \fill[fmGrayBg, draw=fmGray] (-4,\yy-0.18) rectangle (4,\yy+0.18);
    \draw[fmBlue, thick, domain=-4:4, samples=80, smooth] plot (\x, {\yy+0.75*sin(\m*22.5*\x)});
    \draw[fmBlue, thick, dashed, domain=-4:4, samples=80, smooth] plot (\x, {\yy-0.75*sin(\m*22.5*\x)});
    \fill[pattern=north east lines, pattern color=atomLine] (-0.12,\yy+0.18) rectangle (0.12,\yy+0.55);
    \fill[pattern=north east lines, pattern color=atomLine] (-0.12,\yy-0.55) rectangle (0.12,\yy-0.18);
    \node[font=\small, anchor=west] at (4.25,\yy) {\lab};}
  \node[font=\scriptsize, fmGray] at (0,0.75) {узел};
  \node[font=\scriptsize, fmGray] at (-4,1.0) {пучность};
  \node[font=\scriptsize, fmGray] at (4,1.0) {пучность};
  \draw[dimline] (-4,-3.25) -- node[below, font=\scriptsize]{$L$} (4,-3.25);
\end{tikzpicture}
\caption{Распределение амплитуды продольных колебаний по длине стержня,
закреплённого посередине, при первом и втором резонансе (сплошная и
штриховая линии~— два крайних положения)}
\label{fig:rod-res}
\end{figure}

\begin{primer}[Модуль Юнга меди по резонансу]
Медный стержень длиной $L=0{,}1$~м ($\rho=8960$~кг/м$^3$) резонирует на
основной частоте $\nu_\text{рез}=18{,}9$~кГц. По~\eqref{eq:E-res}
$E=4\cdot8960\cdot(1{,}89\cdot10^4)^2\cdot0{,}01\approx1{,}28\cdot10^{11}$~Па~$=128$~ГПа.
\end{primer}

\subsection{Упругость чистых металлов и сплавов}\label{ssec:elastic-alloys}

\paragraph{Модуль как периодическая функция атомного номера}
Модуль упругости при комнатной температуре периодически меняется с
атомным номером (рис.~\ref{fig:young-Z}).
\begin{itemize}
\item В третьем периоде (Na, Mg, Al, Si) модуль растёт с номером: растёт
число валентных электронов и уменьшается атомный радиус.
\item В пределах одной группы \emph{простых} металлов (Be, Mg, Ca, Sr, Ba)
модуль падает с ростом номера вместе с увеличением радиуса. Приближённо
$E=k/a^m$, где $a$~— межатомное расстояние, $k$ и $m$~— постоянные.
\item У \emph{переходных} металлов модуль высок: прочную связь создают
$d$-электроны. Максимум~— у элементов с 5--7 $d$-электронами (Os, Ru, Re,
Ir, W, Mo, Fe, Co). Здесь правило для простых металлов не действует: в
группах Fe--Ru--Os и Co--Rh--Ir модуль \emph{растёт} вместе с атомным
радиусом, что теоретически пока не объяснено.
\end{itemize}

\begin{figure}[ht]
\centering
\begin{tikzpicture}
\begin{axis}[width=0.92\textwidth, height=6.2cm, xmin=0, xmax=86, ymode=log,
  ymin=1, ymax=1000, xlabel={атомный номер $Z$}, ylabel={$E$, ГПа},
  xtick={3,11,19,37,55,83}, unbounded coords=jump]
  \addplot[fmBlue, thick, mark=*, mark size=1.2pt] table[x=Z, y=E] {figures/el-young.dat};
  \foreach \z/\e/\s/\a in {3/4.9/Li/north, 4/287/Be/south, 11/10/Na/north, 13/70/Al/south,
      19/3.5/K/north, 24/279/Cr/south, 26/211/Fe/north, 29/130/Cu/south, 37/2.4/Rb/north,
      42/329/Mo/south, 44/447/Ru/south, 47/83/Ag/north, 55/1.7/Cs/north, 74/411/W/south east,
      76/560/Os/south, 79/78/Au/north, 82/16/Pb/north}{
    \edef\tmp{\noexpand\node[font=\noexpand\scriptsize, anchor=\a] at (axis cs:\z,\e) {\s};}\tmp}
  \node[font=\scriptsize, fmGray] at (axis cs:64,11) {лантаниды};
\end{axis}
\end{tikzpicture}
\caption{Модуль Юнга простых веществ при комнатной температуре в
зависимости от атомного номера (округлённые справочные значения для
поликристаллов; линии соединяют элементы одного периода). Минимумы~—
щелочные металлы, максимумы~— переходные металлы с 5--7 $d$-электронами}
\label{fig:young-Z}
\end{figure}

\paragraph{Фазовые превращения и сплавы}
Фазовые превращения заметно меняют модули: например, превращение
$\alpha$-Fe$\to\gamma$-Fe повышает $E$ из-за более плотной упаковки решётки.
Для сплавов действуют следующие правила.
\begin{itemize}
\item \emph{Непрерывные твёрдые растворы} (одинаковые решётки, близкие
валентности и радиусы): модуль меняется с атомной концентрацией линейно
(Cu--Ni, Cu--Pt, Cu--Au) или почти линейно (Ag--Au). Если в растворе есть
переходные металлы, наблюдается заметное отклонение от аддитивности,
кривая выпукла вверх.
\item \emph{Ограниченные твёрдые растворы}: скорость изменения модуля
$\d E/\d c$ тем больше, чем выше валентность $Z$ растворённого элемента
($\d E/\d c\propto Z^2$), и растёт с разницей атомных радиусов
($\propto\Delta r^2$). Факторы, понижающие температуру плавления сплава,
понижают и модуль: понижение температуры солидуса пропорционально
понижению модуля (при сравнении растворов на основе одного компонента).
\item \emph{Упорядочение} усиливает межатомную связь и повышает модули.
\item \emph{Химические соединения и промежуточные фазы}: чем выше их
температура плавления, тем выше модуль.
\item \emph{Двухфазные сплавы}: модуль аддитивен по объёмным долям фаз.
\end{itemize}

\begin{glavnoe}[Упругость~— структурно нечувствительное свойство]
Модуль упругости определяется силами связи и почти не зависит от
размера зерна однофазного сплава и дисперсности фаз гетерогенного.
Небольшая холодная деформация немного понижает модуль (так же, как у
чистых металлов), а большая может изменить его за счёт текстуры. Этим
упругость резко отличается от структурно чувствительных свойств~—
предела текучести, коэрцитивной силы, внутреннего трения.
\end{glavnoe}

\subsection{Ферромагнитная аномалия упругости}\label{ssec:delta-E}

У ферромагнетика при упругом нагружении помимо обычной деформации
возникает дополнительная \emph{магнитострикционная}: напряжение
поворачивает намагниченность доменов, а с ней меняются и размеры
(п.~\ref{ssec:ferro}). Поэтому модуль размагниченного ферромагнетика
понижен:
\begin{equation}\label{eq:delta-E}
E_\text{ферр}=E_\text{норм}-E',
\end{equation}
где $E'>0$~— ферромагнитная составляющая (\term{$\Delta E$-эффект}). Если
намагнитить образец до насыщения, домены поворачиваться уже не могут,
$E'$ исчезает, и модуль становится «нормальным» (рис.~\ref{fig:delta-E},
кривая~1).

При нагреве к точке Кюри намагниченность падает, $E'$ уменьшается и выше
точки Кюри исчезает. Если $E'$ убывает быстрее, чем $E_\text{норм}$, то в
некотором интервале температур модуль размагниченного ферромагнетика при
нагреве даже \emph{растёт}~— аномалия, хорошо заметная у никеля.

\begin{figure}[ht]
\centering
\begin{tikzpicture}
\begin{axis}[width=0.72\textwidth, height=6.0cm, xmin=0, xmax=800, ymin=170, ymax=230,
  xlabel={$T$, К}, ylabel={$E$, ГПа},
  legend style={at={(0.98,0.97)}, anchor=north east, font=\scriptsize},
  x tick label style={/pgf/number format/1000 sep={}}]
  \addplot[fmRed, thick, domain=0:800, samples=2] {222*(1-2.5e-4*x)};
  \addlegendentry{1~— насыщение ($E_\text{норм}$)}
  \addplot[fmAmber, thick, dashed, domain=0:800, samples=161]
    {222*(1-2.5e-4*x) - 0.35*17*sqrt(max(0,1-(x/631)^4))};
  \addlegendentry{2~— сильное поле}
  \addplot[fmGreen, thick, dash dot, domain=0:800, samples=161]
    {222*(1-2.5e-4*x) - 0.7*17*sqrt(max(0,1-(x/631)^4))};
  \addlegendentry{3~— слабое поле}
  \addplot[fmBlue, very thick, domain=0:800, samples=161]
    {222*(1-2.5e-4*x) - 17*sqrt(max(0,1-(x/631)^4))};
  \addlegendentry{4~— размагниченный}
  \draw[fmGray, densely dotted] (axis cs:631,170) -- (axis cs:631,203);
  \node[font=\scriptsize, anchor=south west, fmGray] at (axis cs:635,172) {$T_\text{К}$};
  \draw[{Stealth[length=1.6mm]}-{Stealth[length=1.6mm]}, fmGray] (axis cs:300,188.8) -- (axis cs:300,205.3);
  \node[font=\scriptsize, anchor=west, fmGray] at (axis cs:305,197) {$E'$};
\end{axis}
\end{tikzpicture}
\caption{$\Delta E$-эффект никеля (схема): температурная зависимость модуля
Юнга размагниченного никеля и никеля, намагниченного в полях
возрастающей напряжённости. В поле насыщения модуль меняется
«нормально»; у размагниченного образца ниже точки Кюри
($T_\text{К}\approx631$~К) модуль понижен на $E'$, а перед точкой Кюри
растёт при нагреве}
\label{fig:delta-E}
\end{figure}
